\chapter*{Guida alla Lettura}
\thispagestyle{empty}
\addcontentsline{toc}{chapter}{Guida alla Lettura}

Per facilitare la consultazione, i concetti del testo sono organizzati in riquadri colorati, ciascuno dedicato a un diverso tipo di contenuto. Ecco il significato di ciascuno:

\begin{definizione}{Concetti e grandezze da ricordare}
Introduce una definizione formale: il significato preciso di una grandezza o di un concetto fisico.
\end{definizione}

\begin{teorema}{Risultati dimostrati}
Raccoglie le proposizioni dimostrate a partire dai principi della fisica, con relativa dimostrazione.
\end{teorema}

\begin{osservazione}{Commenti e precisazioni}
Chiarisce un passaggio, segnala un caso particolare o collega l'argomento ad altre parti del testo.
\end{osservazione}

\begin{esempio}{Applicazioni pratiche}
Esempi ed esercizi svolti, utili per vedere la teoria applicata a un caso concreto. Lo stesso stile è usato per il riquadro \textit{Esercizio}.
\end{esempio}

\begin{attenzione}{Errori comuni}
Segnala un'insidia frequente, un'ipotesi facile da dimenticare o un caso in cui l'intuizione può trarre in inganno.
\end{attenzione}

\begin{notastorica}{Contesto storico}
Un aneddoto o un riferimento storico legato alla scoperta del fenomeno descritto.
\end{notastorica}

\begin{formulario}{Riepilogo di fine capitolo}
Raccoglie in forma compatta le formule principali del capitolo, utile per il ripasso rapido.
\end{formulario}
